\subsection{NN architecture}

We use simple Multilayer Perceptrons as likelihood interpolators, implemented
with the \pytorch package. To ensure stable training dynamics as network width
is varied, we adopt the Maximal Update Parametrization
(\mup)~\cite{Yang2021mup}, in which weight matrices are initialised and
scaled such that activation statistics remain well-behaved independent of the
network width. Concretely, hidden-layer weights are initialised with Kaiming
normal initialisation, while the readout layer weights and biases are
initialised to zero.

Each network takes as input the total yields in all analysis bins (unless some
of them have been removed, see Sec.~\ref{sec:data}) and outputs eight values:
predicted means $\bar\Delta_i$ and corresponding standard deviations
$\sigma_i$ for each of the four $\Delta$s defined in
Eqs.~\ref{eq:delta_exp}--\ref{eq:delta_obsA}. The standard deviations are
passed through a softplus activation to enforce positivity.

The network consists of a number of fully connected blocks, each containing a
dense layer followed by an activation function. The number of blocks and the
number of neurons per block vary depending on the analysis considered, as the
optimal architecture depends on the complexity of the likelihood surface being
interpolated. GELU is used as the activation function for most analyses; for
the \slepton analysis, ELU activations and batch normalisation were found to
perform better. The only hyperparameters optimised on a per-analysis basis are
the L2 regularisation strength and the learning rate, for which a grid search
is performed; all other training hyperparameters are fixed across analyses.

\subsection{Preprocessing}
\label{sec:preprocessing}

Both the input bin yields and the output $\Delta$s are preprocessed before
being fed to the network or used as training targets. We apply the following
two sequential transformations:
%
\begin{enumerate}
  \item \textbf{Signed log transform.} To compress the dynamic range of the
    data while preserving sign information and handling negative values, we
    apply
    %
    \begin{align}
      f(x) = \mathrm{sgn}(x)\,\log(1 + |x|)\,.
      \label{eq:log_trafo}
    \end{align}
    %
    This transformation is exactly invertible via
    $f^{-1}(y) = \mathrm{sgn}(y)\,(e^{|y|}-1)$.
  \item \textbf{Standardisation.} The transformed values are standardised to
    zero mean and unit variance, using statistics computed on the training set.
\end{enumerate}
%
The network therefore operates entirely on standardised log-transformed
quantities. At evaluation time, the predicted means $\bar\Delta_i$ are mapped
back to physical values by applying the inverse transformations in reverse
order. The predicted $\sigma_i$ are kept in the preprocessed space where they
are used to compute the heteroscedastic loss.

\subsection{Training}
\label{sec:training}

We train our networks using the AdamW optimiser with a cosine annealing
learning rate schedule. The learning rate is linearly increased from zero to
the peak value over $4\times10^3$ warm-up steps, after which a cosine decay
is applied over the remaining training steps. We train for $2\times10^5$
optimiser steps with a batch size of 2048. The loss function used is the
heteroscedastic loss,
%
\begin{align}
    \mathcal{L}_\mathrm{het}
    = \frac{1}{\text{batch size}}
      \sum_j^{\text{batch size}}
      \frac{1}{4}\sum_{i}^{4}
      \left[
        \frac{\left(\Delta_{ij}^\mathrm{truth} -
        \bar\Delta_{ij}^\mathrm{prediction}\right)^2}
        {2\,\sigma_{ij}^2}
        + \log \sigma_{ij}
      \right],
    \label{eq:het_loss}
\end{align}
%
where $i$ enumerates the four $\Delta$s, $j$ enumerates samples in a batch,
and $\sigma_{ij}$ is the predicted standard deviation for output $i$ of sample
$j$. The first term penalises residuals weighted by the predicted uncertainty,
while the second term prevents the network from trivially inflating $\sigma$.
In the limit of constant $\sigma$, Eq.~\eqref{eq:het_loss} reduces to the
standard MSE loss. The learned $\sigma(x)$ provides an estimate of the
systematic uncertainty of the surrogate, capturing both noise in the training
data and limitations in the network's expressivity. After training, we load
the weights of the best-performing model as measured on the validation set
and save it to disk together with its metadata as an ONNX file. A detailed
description of the metadata is given in App.~\ref{app:onnx-metadata}.
